\section{The finite identity}\label{sec:finite}

In this section and the next, $q$ is real with $0<q<1$; \cref{thm:main} is
extended to complex $q$ at the end of \cref{sec:limit}. For an integer
$n\ge1$ write $x=q^{2n}$ and define
\begin{equation}\label{eq:f0}
f_0(n)=\frac{q\,B_n\,(1-x^4)^3(x^2-qx+q^2)}{2(1+x^4)(x-q)^2(q^2+x^4)},
\end{equation}
and, for each integer $k\ge0$,
\begin{equation}\label{eq:F}
F(n,k)=f_0(n)\Big(\frac qx\Big)^k\,\frac{1+q^3Q^{2k}/x^3}{1+q^3/x^3}\,
\frac{\poch{qx,\ Q/x^2,\ -Q/x^2,\ -Q/x^2}{Q}{k}}
{\poch{q^3/x,\ q^3/x,\ -q^3/x,\ -q^4/x^4}{Q}{k}} .
\end{equation}
No denominator vanishes: $1-q^3Q^j/x=1-q^{3+2j-2n}$ has an odd exponent, and
the other factors are of the form $1+(\text{positive})$. The factor
$\poch{Q/x^2}{Q}{k}=\prod_{j<k}(1-q^{2j+2-4n})$ contains the factor
$1-q^{0}=0$ (from $j=2n-1$) as soon as $k\ge2n$, so $F(n,k)=0$ for $k\ge2n$.
Hence
\[
T_n=\sum_{k=0}^{2n-1}F(n,k)=\sum_{k\ge0}F(n,k)
\]
is a finite sum.

\begin{lemma}\label{lem:telescope}
For every integer $N\ge1$,
\begin{equation}\label{eq:telescope}
\sum_{n=0}^{N-1}a_n=-T_N .
\end{equation}
\end{lemma}

\begin{proof}
Let $\cert(q,x,y)$ be the rational function defined in \cref{app:cert}, and
put
\[
G(n,k)=\cert(q,q^{2n},q^{2k})\,F(n,k).
\]
The identities proved in \cref{app:cert} say exactly that
\begin{equation}\label{eq:wz}
F(n+1,k)-F(n,k)=G(n,k+1)-G(n,k),\qquad G(n,0)=a_n,
\end{equation}
so $(F,G)$ is a Wilf--Zeilberger pair. Two values of $k$ need care. As a
function of $y=q^{2k}$, the certificate $\cert$ has denominator factors
$x^2-y$ and $q^2x^2-y$, which vanish at $k=2n$ and $k=2n+1$; at exactly these
$k$, the numerator of $F(n,k)$ contains the vanishing factor $1-Q^{2n}/x^2$.
To give $G(n,2n)$ and $G(n,2n+1)$ a meaning, treat $x$ as a real variable near
$q^{2n}$, so that both factors vanish to first order as $x\to q^{2n}$. (The
factor $B_n$ is continued in $x$ by
$\poch{a}{r}{z}=\poch{a}{r}{\infty}/\poch{ar^z}{r}{\infty}$, $r^z=e^{z\log r}$;
it stays finite and nonzero.) The ratio of the two vanishing factors has a
finite nonzero limit, and $G(n,k)$ is defined as the limit of the product.
Every other denominator factor of $\cert$ is nonzero for $0<q<1$. Since
\eqref{eq:wz} holds as an identity of analytic functions of $x$ away from these
points, it continues to hold in the limit.

For $k\ge2n+2$ the certificate is finite and $F(n,k)=0$, so $G(n,k)=0$.
Summing \eqref{eq:wz} over $0\le k\le2n+1$, the right side telescopes to
$G(n,2n+2)-G(n,0)=-a_n$, while the left side is $T_{n+1}-T_n$ (recall
$F(n,2n)=F(n,2n+1)=0$). The two values $G(n,2n)$ and $G(n,2n+1)$ cancel inside
the telescoping sum and are never set to zero. Thus $T_{n+1}-T_n=-a_n$ for
$n\ge1$. Finally, direct substitution in \eqref{eq:F} gives
$T_1=F(1,0)+F(1,1)=-(1-q^5)/[(1+q)(1+q^2)^2]=-a_0$ by \eqref{eq:a0}, and
adding the recurrences gives \eqref{eq:telescope}.
\end{proof}

The finite sum $T_N$ is a terminating ${}_8\phi_7$, so Watson's transformation
applies. For an integer $N\ge1$ put $x=q^{2N}$ and define
\begin{align}
C_N&=\frac{q^2\,x\,B_N\,(1-x^4)^3}{2(1+x^4)(x-q)^2(q^2+x^4)},\label{eq:CN}\\
h_N(k)&=C_N\,Q^k\,
\frac{\poch{Q/x^2}{Q}{k}^2\,\poch{-Q/x^2}{Q}{k}}
{\poch{-Q^2/x^4}{Q}{k}\,\poch{q^3/x}{Q}{k}^2}\qquad(0\le k\le2N-1).\label{eq:hN}
\end{align}

\begin{lemma}\label{lem:watson}
For every integer $N\ge1$,
\begin{equation}\label{eq:watsonsum}
\sum_{n=0}^{N-1}a_n=\sum_{k=0}^{2N-1}h_N(k),
\end{equation}
and every term $h_N(k)$ is nonnegative.
\end{lemma}

\begin{proof}
Apply \cref{fact:watson} with $m=2N-1$ and
\[
a=-q^3/x^3,\qquad b=qx,\qquad c=e=-Q/x^2,\qquad d=Q .
\]
Then $Q^{-m}=Q/x^2$. On the left of \cref{fact:watson}, the parameter $a$
cancels against $aQ/d=a$, and $d=Q$ cancels against the factor $\poch{Q}{Q}{k}$
in the definition of ${}_8\phi_7$. The pair $\poch{\pm Q\sqrt a}{Q}{k}/\poch{\pm\sqrt a}{Q}{k}$
equals $(1-aQ^{2k})/(1-a)=(1+q^3Q^{2k}/x^3)/(1+q^3/x^3)$, and the argument is
$a^2Q^{m+2}/(bcde)=q/x$. Comparing with \eqref{eq:F}, the left side is
$T_N/f_0(N)$. On the right, the prefactor is
\[
\begin{aligned}
\frac{\poch{-q^5/x^3,\ q/x}{Q}{m}}{\poch{-q^3/x^3,\ q^3/x}{Q}{m}}
&=\frac{1+q^3Q^m/x^3}{1+q^3/x^3}\cdot\frac{1-q/x}{1-qQ^m/x}
=\frac{1+q/x}{1+q^3/x^3}\cdot\frac{1-q/x}{1-x/q}\\
&=\frac{x^2}{x^2-qx+q^2}\Big(-\frac qx\Big)
=-\frac{qx}{x^2-qx+q^2},
\end{aligned}
\]
because each ratio telescopes to a single quotient, $Q^m=x^2/q^2$, and
$x^3+q^3=(x+q)(x^2-qx+q^2)$.
In the ${}_4\phi_3$, the parameter $d=Q$ again cancels $\poch{Q}{Q}{k}$. Writing
out the remaining parameters $aQ/(bc)=Q/x^2$, $e=-Q/x^2$, $Q^{-m}=Q/x^2$,
$aQ/b=-Q^2/x^4$, $aQ/c=q^3/x$ and $deQ^{-m}/a=q^3/x$,
we find that $-f_0(N)$ times the prefactor times the $k$-th term of the
${}_4\phi_3$ is exactly $h_N(k)$. Combining with \cref{lem:telescope} gives
\eqref{eq:watsonsum}.

For nonnegativity, note that in \eqref{eq:CN} and \eqref{eq:hN} the factors
that could be negative, $\poch{Q/x^2}{Q}{k}$ and $\poch{q^3/x}{Q}{k}$, occur
squared, and the latter is nonzero. Every other factor is positive when
$0<q<1$ (note $1-x^4>0$).
\end{proof}
